\documentclass[12pt, a4paper]{article}

% ============================================================
%  PAQUETES
% ============================================================
\usepackage[utf8]{inputenc}
\usepackage[T1]{fontenc}
\usepackage[spanish]{babel}
\usepackage{geometry}
\usepackage{graphicx}
\usepackage{hyperref}
\usepackage{xcolor}
\usepackage{listings}
\usepackage{booktabs}
\usepackage{array}
\usepackage{tabularx}
\usepackage{longtable}
\usepackage{fancyhdr}
\usepackage{titlesec}
\usepackage{parskip}
\usepackage{microtype}
\usepackage{enumitem}
\usepackage{tcolorbox}
\usepackage{fontawesome5}
\usepackage{mdframed}
\usepackage{amsmath}
\usepackage{multicol}
\usepackage{caption}
\usepackage{float}

% ============================================================
%  MÁRGENES Y GEOMETRÍA
% ============================================================
\geometry{
  top=2.5cm, bottom=2.5cm,
  left=3cm,  right=3cm,
  headheight=15pt
}

% ============================================================
%  COLORES CORPORATIVOS DE NATIVO
% ============================================================
\definecolor{nativoPrimary}{HTML}{3D9A6F}
\definecolor{nativoDark}{HTML}{2C7A53}
\definecolor{nativoLight}{HTML}{52B484}
\definecolor{nativoPale}{HTML}{EEF5EF}
\definecolor{nativoText}{HTML}{162318}
\definecolor{nativoGray}{HTML}{6B8A72}
\definecolor{codeBg}{HTML}{F7FAF7}
\definecolor{codeFrame}{HTML}{C8E6C9}
\definecolor{warningBg}{HTML}{FFF8DC}
\definecolor{warningFrame}{HTML}{F5C842}
\definecolor{infoBg}{HTML}{E8F4F8}
\definecolor{infoFrame}{HTML}{4A9FD4}

% ============================================================
%  ESTILOS DE CÓDIGO (listings)
% ============================================================
\lstdefinestyle{css}{
  language={[LaTeX]TeX},
  basicstyle=\ttfamily\footnotesize\color{nativoText},
  backgroundcolor=\color{codeBg},
  frame=single,
  framerule=1pt,
  rulecolor=\color{codeFrame},
  keywordstyle=\color{nativoPrimary}\bfseries,
  commentstyle=\color{nativoGray}\itshape,
  stringstyle=\color{nativoDark},
  breaklines=true,
  captionpos=b,
  xleftmargin=12pt,
  xrightmargin=8pt,
  aboveskip=10pt,
  belowskip=10pt,
  showstringspaces=false,
  tabsize=2,
  inputencoding=utf8,
  extendedchars=true,
  literate=
    {á}{{\'a}}1 {é}{{\'e}}1 {í}{{\'i}}1 {ó}{{\'o}}1 {ú}{{\'u}}1
    {Á}{{\'A}}1 {É}{{\'E}}1 {Í}{{\'I}}1 {Ó}{{\'O}}1 {Ú}{{\'U}}1
    {ñ}{{\~n}}1 {Ñ}{{\~N}}1
    {¿}{{?`}}1 {¡}{{!`}}1
    {“}{{``}}1 {”}{{''}}1 {‘}{{`}}1 {’}{{'}}1
    {—}{{---}}1 {–}{{--}}1,
  morekeywords={color, background, font-size, font-family, border-radius,
    display, flex, grid, transition, animation, transform, opacity,
    position, padding, margin, width, height, box-shadow, stroke-dasharray}
}

\lstdefinestyle{javascript}{
  language={},
  basicstyle=\ttfamily\footnotesize\color{nativoText},
  backgroundcolor=\color{codeBg},
  frame=single,
  framerule=1pt,
  rulecolor=\color{codeFrame},
  keywordstyle=\color{nativoDark}\bfseries,
  commentstyle=\color{nativoGray}\itshape,
  stringstyle=\color{nativoPrimary},
  numberstyle=\tiny\color{nativoGray},
  numbers=left,
  numbersep=8pt,
  breaklines=true,
  captionpos=b,
  xleftmargin=20pt,
  xrightmargin=8pt,
  aboveskip=10pt,
  belowskip=10pt,
  showstringspaces=false,
  tabsize=2,
  inputencoding=utf8,
  extendedchars=true,
  literate=
    {á}{{\'a}}1 {é}{{\'e}}1 {í}{{\'i}}1 {ó}{{\'o}}1 {ú}{{\'u}}1
    {Á}{{\'A}}1 {É}{{\'E}}1 {Í}{{\'I}}1 {Ó}{{\'O}}1 {Ú}{{\'U}}1
    {ñ}{{\~n}}1 {Ñ}{{\~N}}1
    {¿}{{?`}}1 {¡}{{!`}}1
    {“}{{``}}1 {”}{{''}}1 {‘}{{`}}1 {’}{{'}}1
    {—}{{---}}1 {–}{{--}}1,
  morekeywords={const, let, function, return, if, else, true, false, null,
    undefined, class, new, throw, try, catch}
}

% ============================================================
%  CAJAS DE CONTENIDO ESPECIALES
% ============================================================
\tcbuselibrary{skins, breakable}

\newtcolorbox{principiobox}[1][]{
  enhanced,
  colback=nativoPale,
  colframe=nativoPrimary,
  boxrule=1.5pt,
  arc=8pt,
  title={\textbf{\color{nativoDark}#1}},
  coltitle=nativoDark,
  fonttitle=\bfseries\small,
  top=8pt, bottom=8pt, left=10pt, right=10pt,
  before skip=12pt, after skip=12pt
}

\newtcolorbox{justificacionbox}[1][]{
  enhanced,
  colback=infoBg,
  colframe=infoFrame,
  boxrule=1.5pt,
  arc=8pt,
  title={\textbf{\color{infoFrame}Justificación: #1}},
  coltitle=infoFrame,
  fonttitle=\bfseries\small,
  top=8pt, bottom=8pt, left=10pt, right=10pt,
  before skip=12pt, after skip=12pt
}

\newtcolorbox{warningbox}[1][]{
  enhanced,
  colback=warningBg,
  colframe=warningFrame,
  boxrule=1.5pt,
  arc=8pt,
  title={\textbf{\color{black}#1}},
  coltitle=black,
  fonttitle=\bfseries\small,
  top=8pt, bottom=8pt, left=10pt, right=10pt,
  before skip=12pt, after skip=12pt
}

% ============================================================
%  ENCABEZADOS Y PIES DE PÁGINA
% ============================================================
\pagestyle{fancy}
\fancyhf{}
\fancyhead[L]{\color{nativoGray}\small\textit{Nativo --- Memoria de Diseño IPO}}
\fancyhead[R]{\color{nativoGray}\small\textit{\thepage}}
\fancyfoot[C]{\color{nativoGray}\footnotesize Interacción Persona-Ordenador}
\renewcommand{\headrulewidth}{0.4pt}
\renewcommand{\footrulewidth}{0.4pt}
\renewcommand{\headrule}{\color{nativoPrimary}\hrule width\headwidth height\headrulewidth}

% ============================================================
%  ESTILO DE TÍTULOS
% ============================================================
\titleformat{\section}
  {\Large\bfseries\color{nativoDark}}
  {\thesection.}{0.8em}{}
  [\vspace{-6pt}\textcolor{nativoPrimary}{\rule{\linewidth}{1.5pt}}\vspace{4pt}]

\titleformat{\subsection}
  {\large\bfseries\color{nativoText}}
  {\thesubsection.}{0.6em}{}

\titleformat{\subsubsection}
  {\normalsize\bfseries\color{nativoGray}}
  {\thesubsubsection.}{0.5em}{}

% ============================================================
%  HYPERREF
% ============================================================
\hypersetup{
  colorlinks=true,
  linkcolor=nativoPrimary,
  urlcolor=nativoDark,
  citecolor=nativoGray,
  pdftitle={Nativo --- Memoria de Diseño IPO},
  pdfauthor={Equipo Nativo},
}

% ============================================================
%  METADATOS DEL DOCUMENTO
% ============================================================
\title{%
  \vspace{-1cm}
  {\Huge\textbf{\textcolor{nativoPrimary}{Nativo}}}\\[0.4cm]
  {\Large\textcolor{nativoGray}{Memoria de Diseño e Implementación}}\\[0.2cm]
  {\large\textcolor{nativoText}{Interacción Persona-Ordenador}}
}

\author{
  \textbf{Breixo Callón Santos --- Angel Enriquez Gago}\\[0.2cm]
  \textcolor{nativoGray}{Grado en Ingeniería Informática}
}

\date{\textcolor{nativoGray}{\today}}

% ============================================================
%  INICIO DEL DOCUMENTO
% ============================================================
\begin{document}

\maketitle
\thispagestyle{empty}

\vspace{0.5cm}
\begin{center}
  \large\textit{\textcolor{nativoGray}{``Camina. Genera energía. Transforma comunidades.''}}
\end{center}


\newpage
\tableofcontents
\newpage

% ============================================================
\section{Introducción y Concepto de la Aplicación}
% ============================================================

\subsection{Descripción General}

\textbf{Nativo} es una aplicación móvil (con versión responsiva de escritorio) que convierte la actividad física diaria del usuario en impacto social real. El ecosistema de la aplicación funciona bajo la siguiente cadena de valor:

\begin{center}
\begin{mdframed}[backgroundcolor=nativoPale, linecolor=nativoPrimary, linewidth=1.5pt, roundcorner=8pt]
\begin{center}
\vspace{6pt}
{\Large\textbf{\textcolor{nativoDark}{Pasos \quad$\rightarrow$\quad Energía \quad$\rightarrow$\quad Créditos \quad$\rightarrow$\quad Impacto Social}}}
\vspace{6pt}
\end{center}
\end{mdframed}
\end{center}

\subsection{Propuesta de Valor}

La propuesta de valor de Nativo se sustenta en tres pilares fundamentales identificados durante el proceso de \textit{needfinding}:

\begin{enumerate}[leftmargin=*, label=\textbf{\arabic*.}]
  \item \textbf{Motivación intrínseca}: El usuario ve el resultado de su actividad física de forma inmediata (créditos generados), lo que refuerza el comportamiento positivo.
  \item \textbf{Impacto tangible}: Los créditos se destinan a proyectos comunitarios reales, cerrando el ciclo de significado entre el esfuerzo individual y el beneficio colectivo.
  \item \textbf{Gamificación ética}: A diferencia de apps de fitness convencionales, Nativo añade una dimensión social que transforma el ejercicio en altruismo.
\end{enumerate}

\subsection{Conversión de Unidades}

La mecánica central de Nativo opera con la siguiente fórmula de conversión:

\begin{equation}
  C = \left\lfloor \frac{P}{100} \right\rfloor
\end{equation}

Donde $C$ es el número de créditos Nativo y $P$ es el número de pasos totales acumulados. Esta conversión es deliberadamente simple para que el usuario la interiorice sin esfuerzo cognitivo (100 pasos = 1 crédito).

%=============================
\newpage

\section{Needfinding}

La fase de \textit{needfinding} tiene como objetivo identificar las necesidades, expectativas, motivaciones y posibles dificultades de los usuarios potenciales de Nativo. En concreto, este estudio pretende:

\begin{itemize}
    \item Analizar el interés de los usuarios por aplicaciones relacionadas con actividad física.
    \item Evaluar el atractivo de un sistema de recompensas basado en créditos.
    \item Comprobar si el impacto social de la aplicación constituye un factor motivador.
    \item Detectar posibles problemas de comprensión en la relación entre pasos, energía, créditos y proyectos.
    \item Obtener información útil para justificar futuras decisiones de diseño centrado en el usuario.
\end{itemize}

Se plantea una muestra de \textbf{30 usuarios}, divididos en tres perfiles distintos.

\begin{itemize}
    \item \textbf{Estudiantes}: usuarios jóvenes familiarizados con aplicaciones móviles y con interés moderado en herramientas digitales de organización y seguimiento.
    \item \textbf{Personas activas o deportistas}: usuarios que ya utilizan aplicaciones de salud o actividad física y están habituados a métricas y objetivos.
    \item \textbf{Usuarios casuales}: personas que no usan este tipo de aplicaciones de forma habitual y cuya motivación inicial es menor.
\end{itemize}

Cuestionario utilizado en la muestra.

\subsubsection{Bloque 1: hábitos de uso}

\begin{enumerate}[label=\arabic*.]
    \item ¿Utilizas aplicaciones para medir tu actividad física diaria?
    \item ¿Con qué frecuencia consultas este tipo de aplicaciones?
    \item ¿Qué valoras más en una aplicación de actividad: simplicidad, estadísticas, recompensas o diseño visual?
\end{enumerate}

\subsubsection{Bloque 2: motivación y recompensa}

\begin{enumerate}[label=\arabic*.]
    \setcounter{enumi}{3}
    \item ¿Te sentirías más motivado/a a caminar o moverte si obtuvieras recompensas digitales?
    \item ¿Qué tipo de recompensa te resulta más atractiva?
    \begin{itemize}
        \item Recompensa económica
        \item Recompensa simbólica o logros
        \item Impacto social en la comunidad
    \end{itemize}
    \item ¿Consideras importante poder visualizar tu progreso diario de manera inmediata?
\end{enumerate}

\subsubsection{Bloque 3: impacto social}

\begin{enumerate}[label=\arabic*.]
    \setcounter{enumi}{6}
    \item ¿Te interesa que tu actividad diaria contribuya a proyectos comunitarios?
    \item ¿Preferirías elegir manualmente a qué proyecto donar tus créditos o que el sistema lo haga automáticamente?
    \item ¿Te parece importante conocer exactamente cuántos créditos has aportado y a qué proyecto han ido destinados?
\end{enumerate}

\subsubsection{Bloque 4: comprensión y usabilidad}

\begin{enumerate}[label=\arabic*.]
    \setcounter{enumi}{9}
    \item ¿Te parece fácil de entender la relación entre pasos, energía, créditos y donaciones?
    \item ¿Qué elemento te ayudaría más a comprender el sistema: una explicación inicial, iconos claros, barras de progreso o ejemplos visuales?
    \item ¿Qué tan importante es para ti que la interfaz sea simple y rápida de usar?
\end{enumerate}


\vspace{0.3cm}
\noindent\textcolor{nativoPrimary}{\rule{\linewidth}{1.5pt}}
\vspace{0.3cm}

\begin{table}[h!]
\centering
\begin{tabular}{p{8cm}c}
\toprule
\textbf{Factor evaluado} & \textbf{Media (1-5)} \\
\midrule
Importancia de ver impacto real en la comunidad & 4,6 \\
Motivación por obtener recompensas & 4,2 \\
Facilidad de uso de la aplicación & 4,8 \\
Interés en donar créditos a proyectos & 3,9 \\
Comprensión inicial del sistema de créditos & 3,5 \\
\bottomrule
\end{tabular}
\caption{Valoraciones medias simuladas de los usuarios.}
\end{table}

% ============================
\newpage
% ============================================================
\section{Principios de Diseño Aplicados}
% ============================================================

\subsection{Marco Teórico: Diseño Centrado en el Usuario (DCU)}

La aplicación fue diseñada siguiendo los principios del DCU según la norma \textbf{ISO 9241-210:2010}. Esto implica que cada decisión de diseño se toma teniendo en cuenta las necesidades, capacidades y limitaciones de los usuarios finales.

\begin{principiobox}[Principios de Nielsen aplicados a Nativo]
\begin{enumerate}[leftmargin=*, label=\textbf{N\arabic*.}, itemsep=4pt]
  \item \textbf{Visibilidad del estado del sistema}: El anillo de progreso y los créditos se actualizan en tiempo real.
  \item \textbf{Consistencia y estándares}: Navegación inferior tipo iOS/Android, patrones de interacción familiares.
  \item \textbf{Reconocimiento antes que recuerdo}: Los iconos pictóricos reducen la carga cognitiva.
  \item \textbf{Diseño estético y minimalista}: Cada elemento visual tiene una función específica.
  \item \textbf{Feedback adecuado}: Toasts, animaciones y microinteracciones informan al usuario inmediatamente.
\end{enumerate}
\end{principiobox}

\subsection{Modelo Mental del Usuario}

Basándonos en el needfinding, el modelo mental del usuario de Nativo es:

\begin{itemize}
  \item El movimiento físico tiene un valor que va más allá de la salud personal.
  \item Las pequeñas acciones individuales pueden acumularse en cambios sociales significativos.
  \item Las aplicaciones de fitness deben proporcionar recompensas tangibles e inmediatas.
\end{itemize}

\newpage
% ============================================================
\section{Justificación Tipográfica}
% ============================================================

\subsection{Selección de Tipografías}

Nativo utiliza un sistema tipográfico de dos fuentes con roles claramente diferenciados:

\begin{table}[H]
\centering
\caption{Sistema tipográfico de Nativo}
\begin{tabularx}{\textwidth}{>{\bfseries}l X X}
\toprule
\textbf{Fuente} & \textbf{Rol} & \textbf{Justificación} \\
\midrule
\textbf{Sora} & Tipografía principal (UI, textos, etiquetas) & Humanista geométrica con personalidad propia; equilibrio entre calidez y modernidad \\
\addlinespace
\textbf{JetBrains Mono} & Datos numéricos (pasos, créditos, porcentajes) & Alta legibilidad para cifras; diferenciación semántica del dato numérico \\
\bottomrule
\end{tabularx}
\end{table}

\subsection{Justificación de Sora como Tipografía Principal}

\begin{justificacionbox}[Sora como fuente principal]
\textbf{Sora} fue elegida sobre alternativas más comunes (Inter, Roboto, SF Pro) por las siguientes razones:

\begin{itemize}
  \item \textbf{Carácter único}: Sus terminaciones redondeadas evocan organicidad, alineándose con la temática de sostenibilidad de la aplicación.
  \item \textbf{Legibilidad en pantallas pequeñas}: Diseñada específicamente para entornos digitales, con métricas optimizadas para tamaños pequeños.
  \item \textbf{Rango de pesos}: De 300 a 800, permite establecer jerarquías visuales claras sin cambiar de fuente.
  \item \textbf{Diferenciación}: Evita la genericidad de Inter (sobreutilizada en apps).
  \item \textbf{Congruencia de marca}: Las formas orgánicas de Sora son coherentes con los valores de la marca (naturaleza, comunidad, movimiento).
\end{itemize}
\end{justificacionbox}

\subsection{Justificación de JetBrains Mono para Datos Numéricos}

\begin{justificacionbox}[JetBrains Mono para cifras]
\begin{itemize}
  \item \textbf{Monoespaciado}: Evita el \textit{shifting} visual cuando los números cambian (el ancho del dígito es constante), lo que hace que las animaciones numéricas sean más suaves y profesionales.
  \item \textbf{Diferenciación semántica}: El contraste entre Sora (humanista) y JetBrains Mono (técnica) comunica visualmente la diferencia entre \textit{contenido narrativo} y \textit{dato cuantitativo}.
  \item \textbf{Legibilidad de cifras}: Distingue claramente entre 0, O, 1, l, I —crucial para lecturas rápidas de pasos y créditos.
  \item \textbf{Peso visual equilibrado}: A tamaño equivalente, JetBrains Mono tiene mayor peso visual, lo que hace que los números ``destaquen'' naturalmente.
\end{itemize}
\end{justificacionbox}

\subsection{Escala Tipográfica}

\begin{table}[H]
\centering
\caption{Escala tipográfica utilizada en Nativo}
\begin{tabular}{llll}
\toprule
\textbf{Uso} & \textbf{Fuente} & \textbf{Tamaño} & \textbf{Peso} \\
\midrule
Número de pasos (anillo) & JetBrains Mono & 2.2rem & 800 \\
Créditos disponibles     & JetBrains Mono & 2.4rem & 800 \\
Título de sección        & Sora           & 1.5rem & 700 \\
Nombre de proyecto       & Sora           & 1.0rem & 700 \\
Descripción de proyecto  & Sora           & 0.78rem & 400 \\
Etiqueta de navegación   & Sora           & 0.62rem & 600 \\
Estadísticas secundarias & JetBrains Mono & 1.1rem & 700 \\
\bottomrule
\end{tabular}
\end{table}

\newpage
% ============================================================
\section{Justificación del Color}
% ============================================================

\subsection{Paleta Principal}

\begin{table}[H]
\centering
\caption{Paleta de colores de Nativo}
\begin{tabularx}{\textwidth}{lllX}
\toprule
\textbf{Rol} & \textbf{Nombre} & \textbf{Hex} & \textbf{Uso} \\
\midrule
\colorbox{nativoPrimary}{\color{white}\texttt{\#3D9A6F}} & Verde primario & \texttt{\#3D9A6F} & Botones, anillo, acentos principales \\
\colorbox{nativoDark}{\color{white}\texttt{\#2C7A53}} & Verde oscuro & \texttt{\#2C7A53} & Texto sobre fondo claro, degradados \\
\colorbox{nativoLight}{\color{white}\texttt{\#52B484}} & Verde claro & \texttt{\#52B484} & Estados hover, degradados \\
\colorbox{nativoPale}{\color{nativoText}\texttt{\#EEF5EF}} & Verde pálido & \texttt{\#EEF5EF} & Fondo de la aplicación \\
\colorbox[HTML]{F5C842}{\color{black}\texttt{\#F5C842}} & Dorado & \texttt{\#F5C842} & Racha diaria, animaciones de crédito \\
\colorbox[HTML]{4A9FD4}{\color{white}\texttt{\#4A9FD4}} & Azul & \texttt{\#4A9FD4} & Proyectos de educación, toasts info \\
\bottomrule
\end{tabularx}
\end{table}

\subsection{Fundamentos Psicológicos del Color Verde}

\begin{justificacionbox}[El verde como color de marca]
\textbf{Psicología del color} (Fehrman \& Fehrman, 2004; Elliot \& Maier, 2012):

\begin{itemize}
  \item \textbf{Naturaleza y sostenibilidad}: El verde es el color universalmente asociado al medio ambiente. Para una app de impacto ecológico, esta asociación es congruente y reduce la disonancia cognitiva.
  \item \textbf{Vitalidad y salud}: Evoca movimiento, actividad física y bienestar, alineándose con el contador de pasos.
  \item \textbf{Confianza y crecimiento}: Transmite solidez y progreso, reforzando la sensación de que los créditos se acumulan de forma segura.
  \item \textbf{Accesibilidad}: El verde \texttt{\#3D9A6F} sobre blanco cumple el ratio de contraste WCAG AA (4.5:1) para texto estándar.
\end{itemize}

\textbf{Por qué no el verde de \texttt{\#5cb85c} original}:
El verde \texttt{\#3D9A6F} fue preferido sobre el verde bootstrap \texttt{\#5cb85c} porque es menos genérico y tiene mayor sofisticación cromática, evitando la asociación con templates prediseñados.
\end{justificacionbox}

\subsection{Uso del Degradado en el Panel de Créditos}

El panel de créditos utiliza un degradado:

\begin{center}
\texttt{linear-gradient(135deg, \#2C7A53 -> \#3D9A6F -> \#52B484)}
\end{center}

Esta decisión se justifica por:
\begin{enumerate}
  \item \textbf{Jerarquía visual}: El elemento más importante de la pantalla recibe el tratamiento visual más elaborado.
  \item \textbf{Profundidad}: El degradado crea sensación de tridimensionalidad sin requerir sombras excesivas.
  \item \textbf{Energía}: El flujo de color de izquierda a derecha evoca movimiento y progreso.
\end{enumerate}

\subsection{Color para Proyectos Comunitarios}

Cada proyecto utiliza un color temático distinto, aportando:
\begin{itemize}
  \item \textbf{Diferenciación}: El usuario identifica proyectos por color, reduciendo la carga cognitiva.
  \item \textbf{Significado semántico}: Verde (alimentación/naturaleza), azul (educación/agua), naranja (energía solar), púrpura (salud/agua limpia).
  \item \textbf{Consistencia}: Cada color se aplica tanto al icono, la barra de progreso y el botón del proyecto, creando cohesión.
\end{itemize}

\newpage
% ============================================================
\section{Justificación CSS: Técnicas y Decisiones}
% ============================================================

\subsection{Variables CSS Custom Properties}

Nativo utiliza un sistema de tokens de diseño implementado con variables CSS:

\begin{lstlisting}[style=css, caption={Variables CSS de Nativo (extracto)}]
:root {
  --primary:        #3D9A6F;
  --primary-dark:   #2C7A53;
  --primary-glow:   rgba(61, 154, 111, 0.25);
  --font:           'Sora', -apple-system, sans-serif;
  --mono:           'JetBrains Mono', monospace;
  --ring-circumf:   565.49; /* 2 * PI * 90 */
  --ease-bounce:    cubic-bezier(0.34, 1.56, 0.64, 1);
  --t-mid:          0.3s var(--ease-out);
}
\end{lstlisting}

\begin{justificacionbox}[Uso de Custom Properties]
\begin{itemize}
  \item \textbf{Mantenibilidad}: Cambiar el color primario en un solo lugar actualiza toda la UI.
  \item \textbf{Rendimiento}: Las variables CSS son procesadas por el motor de renderizado del navegador de forma más eficiente que sus equivalentes preprocesados.
  \item \textbf{Theming}: Facilita la implementación futura de modo oscuro con una simple sustitución de variables.
  \item \textbf{Semántica}: Los nombres \texttt{--primary}, \texttt{--surface}, \texttt{--text-1} comunican el rol del token, no su valor concreto.
\end{itemize}
\end{justificacionbox}

\subsection{Anillo de Progreso SVG}

La técnica del anillo utiliza \texttt{stroke-dasharray} y \texttt{stroke-dashoffset}:

\begin{lstlisting}[style=css, caption={Técnica SVG para el anillo de progreso}]
/* Circunferencia = 2 * PI * radio = 2 * 3.14159 * 90 = 565.49 */
.ring-fill {
  stroke-dasharray:  565.49;
  stroke-dashoffset: 565.49; /* 100% = vacío */
  transition: stroke-dashoffset 0.8s cubic-bezier(0.22, 1, 0.36, 1);
}
/* Para P% de progreso:
   offset = circumferencia * (1 - P/100) */
\end{lstlisting}

\begin{justificacionbox}[Anillo SVG vs. Canvas vs. CSS puro]
\textbf{¿Por qué SVG?}
\begin{itemize}
  \item \textbf{Escalabilidad}: SVG escala perfectamente a cualquier resolución sin pérdida de calidad.
  \item \textbf{Animación CSS}: \texttt{stroke-dashoffset} es animable con CSS, aprovechando la aceleración por hardware del navegador (compositor thread).
  \item \textbf{Accesibilidad}: Los elementos SVG admiten atributos ARIA (\texttt{aria-label}) para lectores de pantalla.
  \item \textbf{Gradiente}: Los degradados en SVG (\texttt{linearGradient}) permiten efectos que el CSS \texttt{border-image} no puede conseguir con la misma fidelidad en bordes circulares.
\end{itemize}
\end{justificacionbox}

\subsection{Layout Responsivo: Mobile-First}

La aplicación usa un enfoque \textit{mobile-first}:

\begin{lstlisting}[style=css, caption={Estructura responsiva de Nativo}]
/* Base: móvil (pantalla completa) */
.phone-mockup {
  width: 100%;
  background: none;
  border-radius: 0;
}

/* A partir de 768px: escritorio con mockup */
@media (min-width: 768px) {
  .desktop-sidebar { display: flex; }
  .phone-mockup {
    width: 390px;
    background: #1A1A1A;
    border-radius: 52px;
  }
}
\end{lstlisting}

\begin{justificacionbox}[Estrategia Mobile-First]
El 87\% de los usuarios de apps de fitness acceden desde dispositivos móviles (\href{https://www.statista.com/outlook/hmo/digital-health/digital-fitness-well-being/health-wellness-coaching/fitness-apps/worldwide}{Statista, 2024}). El diseño mobile-first garantiza que la experiencia principal esté optimizada para el contexto de uso real: el usuario caminando con su teléfono.

La vista de escritorio añade valor informativo (sidebar con propuesta de valor) sin comprometer la experiencia móvil.
\end{justificacionbox}

\subsection{Animaciones y Transiciones}

\begin{table}[H]
\centering
\caption{Animaciones implementadas y su justificación UX}
\begin{tabularx}{\textwidth}{llX}
\toprule
\textbf{Elemento} & \textbf{Técnica CSS} & \textbf{Propósito UX} \\
\midrule
Anillo de progreso & \texttt{stroke-dashoffset} transition & Feedback inmediato del progreso de pasos \\
Número de créditos & \texttt{@keyframes} scale bounce & Refuerzo positivo al ganar créditos \\
Partículas de donación & \texttt{@keyframes} translate + fade & Celebración del acto altruista \\
Cambio de tab & \texttt{@keyframes} slideInUp & Orientación espacial entre pantallas \\
Bottom sheet modal & \texttt{transform translateY} & Patrón nativo de iOS/Android \\
Planeta orbitando & \texttt{@keyframes} translateY float & Metáfora visual de impacto global \\
Órbitas animadas & \texttt{@keyframes} rotate & Énfasis en la naturaleza sistémica del impacto \\
\bottomrule
\end{tabularx}
\end{table}

\begin{lstlisting}[style=css, caption={Curvas de easing custom de Nativo}]
:root {
  /* Elástico: ideal para microinteracciones táctiles */
  --ease-bounce: cubic-bezier(0.34, 1.56, 0.64, 1);
  /* Desaceleración suave: ideal para transiciones de pantalla */
  --ease-out: cubic-bezier(0.22, 1, 0.36, 1);
}
\end{lstlisting}

\begin{justificacionbox}[Por qué curvas de easing custom]
Las curvas predefinidas (\texttt{ease}, \texttt{ease-in-out}) resultan genéricas y «muertas». Las curvas custom:
\begin{itemize}
  \item \texttt{cubic-bezier(0.34, 1.56, 0.64, 1)}: Tiene un ligero \textit{overshoot} que simula la física de objetos con masa. Comunica respuesta táctil en botones y elementos interactivos.
  \item \texttt{cubic-bezier(0.22, 1, 0.36, 1)}: Arranca rápido y decelera suavemente, ideal para animaciones de entrada que deben sentirse «naturales» y no mecánicas.
\end{itemize}
\end{justificacionbox}

\subsection{Bottom Sheet como Patrón Modal}

El modal de donación implementa el patrón \textit{Bottom Sheet}:

\begin{lstlisting}[style=css, caption={Bottom Sheet: técnica CSS}]
.modal-sheet {
  border-radius: 28px 28px 0 0;
  transform: translateY(100%);        /* Oculto debajo */
  transition: transform 0.4s var(--ease-out);
}

.modal-backdrop.open .modal-sheet {
  transform: translateY(0);           /* Slide up al abrir */
}
\end{lstlisting}

\begin{justificacionbox}[Bottom Sheet vs. Modal centrado]
El \textit{Bottom Sheet} fue elegido sobre un modal centrado por:
\begin{itemize}
  \item \textbf{Alcance con el pulgar}: Los estudios de ergonomía móvil (Hoober, 2013) demuestran que la zona inferior de la pantalla es la más accesible para interacción con una sola mano.
  \item \textbf{Contexto preservado}: Al mostrar parcialmente el contenido detrás, el usuario no pierde el contexto de qué proyecto está evaluando.
  \item \textbf{Convención nativa}: Es el patrón estándar en iOS (Sheet) y Material Design (Bottom Sheet), por lo que los usuarios lo reconocen sin aprendizaje.
\end{itemize}
\end{justificacionbox}

\newpage
% ============================================================
\section{Justificación JavaScript: Arquitectura y Lógica}
% ============================================================

\subsection{Patrón de Estado Centralizado}

Nativo implementa un patrón de \textbf{estado centralizado} sin frameworks:

\begin{lstlisting}[style=javascript, caption={Estado centralizado de Nativo}]
const state = {
  steps: 0,             // Pasos acumulados
  stepsGoal: 10000,     // Objetivo diario
  credits: 50,          // Créditos disponibles
  creditsEarnedToday: 0,
  STEPS_PER_CREDIT: 100,
  projects: [ /* ... */ ],
  modal: { open: false, projectId: null, selectedAmount: 0 },
  donationHistory: [],
  activeTab: 'home',
};
\end{lstlisting}

\begin{justificacionbox}[Estado centralizado vs. estado distribuido]
\begin{itemize}
  \item \textbf{Predecibilidad}: La UI es una función del estado. Dado un estado, la interfaz es siempre la misma. Esto facilita el debugging.
  \item \textbf{Consistencia}: Evita inconsistencias entre el número de créditos mostrado en distintas pantallas.
  \item \textbf{Testabilidad}: El estado puede exportarse y compararse antes/después de cada operación.
  \item \textbf{Escalabilidad}: La adición de nuevas funcionalidades solo requiere extender el objeto \texttt{state}.
\end{itemize}
\end{justificacionbox}

\subsection{Función \texttt{addSteps(n)}}

\begin{lstlisting}[style=javascript, caption={Implementación de addSteps}]
function addSteps(n) {
  const prevSteps = state.steps;
  state.steps      += n;
  state.stepsToday += n;

  // Créditos nuevos = diferencia entre total posible antes y despues
  const totalPossible = Math.floor(state.steps / state.STEPS_PER_CREDIT);
  const prevPossible  = Math.floor(prevSteps   / state.STEPS_PER_CREDIT);
  const newCredits    = totalPossible - prevPossible;

  if (newCredits > 0) {
    state.credits            += newCredits;
    state.creditsEarnedToday += newCredits;
    showCreditFloat('+' + newCredits + ' coin');
    showToast(`+${newCredits} creditos generados`, 'success');
  }
  // Detectar hitos (2500, 5000, 7500, 10000)
  milestones.forEach(m => {
    if (prevSteps < m && state.steps >= m) triggerMilestone(m);
  });
  // Actualizar toda la UI
  updateRing(); updateEnergy(); updateCreditsPanel();
}
\end{lstlisting}

\begin{justificacionbox}[Cálculo de créditos por diferencia]
El cálculo de créditos usa \textbf{diferencia de mínimos enteros} en lugar de dividir directamente los pasos añadidos. Esto garantiza exactitud cuando los pasos se acumulan en lotes irregulares: si el usuario tiene 95 pasos y añade 10, no gana crédito aún, pero si añade 10 más llega a 115 y gana exactamente 1.
\end{justificacionbox}

\subsection{Función \texttt{donate(projectId, amount)}}

\begin{lstlisting}[style=javascript, caption={Lógica de donación}]
function confirmDonate() {
  const amount  = state.modal.selectedAmount;
  const project = state.projects.find(p => p.id === state.modal.projectId);

  // Validaciones defensivas
  if (amount > state.credits) {
    showToast('No tienes suficientes creditos', 'error');
    applyShakeAnimation(DOM.creditsDisp());
    return;
  }

  // Ajuste: no superar la meta del proyecto
  const donateAmt = Math.min(amount, project.goal - project.funded);

  // Mutacion del estado
  state.credits        -= donateAmt;
  state.creditsDonated += donateAmt;
  project.funded       += donateAmt;

  // Registrar en historial
  state.donationHistory.push({ /* ... metadatos ... */ });

  // Efectos visuales y feedback
  triggerDonationParticles();
  // [Actualizar UI completa]
}
\end{lstlisting}

\begin{justificacionbox}[Validaciones y ajuste de donación]
\begin{itemize}
  \item \textbf{Validación de saldo}: Previene estados inconsistentes (créditos negativos).
  \item \textbf{Ajuste de exceso}: \texttt{Math.min(amount, meta - financiado)} garantiza que ningún proyecto sea sobrefinanciado, manteniendo la integridad del modelo de datos.
  \item \textbf{Animación de error}: El efecto \textit{shake} ante fondos insuficientes proporciona feedback háptico visual, más informativo que un mensaje de texto simple.
\end{itemize}
\end{justificacionbox}

\subsection{Función \texttt{setTab(tabId)}}

\begin{lstlisting}[style=javascript, caption={Navegación entre pestañas}]
function setTab(tabId) {
  if (state.activeTab === tabId) return; // Evitar re-render innecesario

  state.activeTab = tabId;

  // Actualizar pantallas visibles
  DOM.tabScreens().forEach(screen => {
    screen.classList.toggle('active', screen.id === `tab-${tabId}`);
  });

  // Actualizar indicadores de navegacion
  DOM.navBtns().forEach(btn => {
    btn.classList.toggle('active', btn.dataset.tab === tabId);
  });

  // Re-renderizar contenido especifico del tab activo
  if (tabId === 'projects') renderProjects();
  if (tabId === 'impact')   { updateImpact(); renderDonationHistory(); }
}
\end{lstlisting}

\begin{justificacionbox}[Renderizado bajo demanda]
El patrón \textit{render-on-navigate} solo actualiza el contenido del tab activo cuando se accede a él, en lugar de mantener todos los tabs sincronizados permanentemente. Esto reduce el coste computacional y evita renders innecesarios, lo que es especialmente importante en dispositivos móviles de gama media-baja.
\end{justificacionbox}

\subsection{Sistema de Toasts (Notificaciones)}

\begin{lstlisting}[style=javascript, caption={Sistema de toasts dinámicos}]
function showToast(message, type = 'success', duration = 3000) {
  const toast = document.createElement('div');
  toast.className = `toast toast-${type}`;
  toast.innerHTML = `<span>${icons[type]}</span><span>${message}</span>`;

  DOM.toastContainer().appendChild(toast);

  setTimeout(() => {
    toast.classList.add('hiding');
    setTimeout(() => toast.remove(), 300);
  }, duration);
}
\end{lstlisting}

\begin{justificacionbox}[Feedback inmediato: el rol de los toasts]
Los toasts cumplen el principio de Nielsen N1 (\textit{Visibilidad del estado del sistema}) proporcionando:
\begin{itemize}
  \item \textbf{Confirmación de acción}: El usuario sabe que su donación fue procesada.
  \item \textbf{Gamificación}: El toast de créditos ganados (+N créditos) actúa como refuerzo positivo inmediato.
  \item \textbf{Error handling amigable}: Los mensajes de error son concisos y explican la causa.
  \item \textbf{No intrusivos}: Los toasts no bloquean la interfaz; aparecen y desaparecen automáticamente.
\end{itemize}
\end{justificacionbox}

\newpage
% ============================================================
\section{Arquitectura de la Información}
% ============================================================

\subsection{Estructura de Navegación}

La arquitectura de información de Nativo sigue un modelo \textbf{Tab Navigation} de 3 nodos:

\begin{center}
\begin{mdframed}[backgroundcolor=nativoPale, linecolor=nativoPrimary, linewidth=1pt, roundcorner=6pt]
\begin{center}
\vspace{8pt}
\textbf{Inicio} (\faHome) \quad $\leftrightarrow$ \quad \textbf{Proyectos} (\faLeaf) \quad $\leftrightarrow$ \quad \textbf{Impacto} (\faChartBar)
\vspace{4pt}

\textit{Pasos + Créditos} \quad\quad\quad \textit{Donación} \quad\quad\quad \textit{Historial}
\vspace{8pt}
\end{center}
\end{mdframed}
\end{center}

\subsection{Jerarquía Visual por Pantalla}

\subsubsection{Pantalla Inicio}

\begin{enumerate}[label=\textbf{\arabic*.}]
  \item Anillo de progreso (elemento focal, mayor tamaño visual)
  \item Panel de créditos (segundo en jerarquía, degradado que atrae la atención)
  \item Botones de pasos (acción primaria, nivel inferior de jerarquía)
\end{enumerate}

Esta jerarquía responde al flujo de atención natural del usuario: primero quiere saber su progreso, luego su recompensa, y después ejecutar una acción.

\subsubsection{Pantalla Proyectos}

Las tarjetas de proyectos están diseñadas siguiendo el patrón \textit{Card UI} con tres zonas cognitivas:
\begin{itemize}
  \item \textbf{Zona de identidad}: Icono + nombre + descripción + etiqueta temática
  \item \textbf{Zona de progreso}: Barra de progreso + cifras de financiación
  \item \textbf{Zona de acción}: Porcentaje financiado + botón de donación
\end{itemize}

\newpage
% ============================================================
\section{Evaluación de Usabilidad}
% ============================================================

\subsection{Métricas de Diseño Evaluadas}

\begin{table}[H]
\centering
\caption{Checklist de principios de usabilidad implementados}
\begin{tabularx}{\textwidth}{>{\bfseries}lXc}
\toprule
\textbf{Principio} & \textbf{Implementación en Nativo} & \textbf{Estado} \\
\midrule
Feedback inmediato & Anillo, toasts, créditos flotantes, bounce animation & \faCheckCircle \\
Consistencia & Sistema de tokens CSS, paleta unificada, tipografía dual & \faCheckCircle \\
Minimizar carga cognitiva & Iconos pictóricos, etiquetas cortas, jerarquía clara & \faCheckCircle \\
Control del usuario & Botón cancelar en modal, confirmación antes de donar & \faCheckCircle \\
Prevención de errores & Botón deshabilitado si monto = 0 o > saldo disponible & \faCheckCircle \\
Accesibilidad básica & ARIA labels, contraste WCAG AA, \texttt{aria-live} en toasts & \faCheckCircle \\
Gestos nativos & Swipe down para cerrar modal (iOS behavior) & \faCheckCircle \\
Estados vacíos & Pantalla de impacto con CTA cuando no hay donaciones & \faCheckCircle \\
Gamificación ética & Rachas, logros, partículas de celebración & \faCheckCircle \\
Patrones reconocibles & Bottom sheet, tab bar, cards: convenciones establecidas & \faCheckCircle \\
\bottomrule
\end{tabularx}
\end{table}

\subsection{Decisiones de Microinteracción}

\begin{warningbox}[Las microinteracciones como diferencial de calidad]
Las microinteracciones no son decorativas: son comunicación funcional. En Nativo, cada microinteracción tiene una justificación UX específica:

\begin{itemize}
  \item \textbf{Scale bounce en créditos}: El número «rebota» al ganar créditos, simulando un objeto físico que aterriza. Esta metáfora física comunica peso y valor.
  \item \textbf{Crédito flotante (+N créditos)}: Aparece en la posición exacta del número de créditos y asciende, conectando visualmente la causa (pasos) con el efecto (crédito).
  \item \textbf{Partículas de donación}: 8 emojis salen radialmente del botón de donar, evocando la «dispersión» del impacto hacia la comunidad. Es una metáfora visual de la descentralización del bien.
  \item \textbf{Milestone pulse en el anillo}: El anillo pulsa luminosamente al alcanzar un hito (2.500, 5.000, 7.500, 10.000 pasos), diferenciando un evento especial de un progreso ordinario.
\end{itemize}
\end{warningbox}

\newpage
% ============================================================
\section{Conclusiones}
% ============================================================

Nativo es un ejemplo de cómo el Diseño Centrado en el Usuario puede convertirse en el motor de adopción de comportamientos pro-sociales. Cada decisión de diseño —desde la elección tipográfica hasta la curva de easing de una animación— está al servicio de un objetivo claro: hacer que caminar sea \textit{significativo}.

Las decisiones más críticas pueden resumirse en:

\begin{enumerate}[leftmargin=*, itemsep=6pt]
  \item \textbf{Tipografía dual Sora + JetBrains Mono}: Diferencia contenido narrativo de dato cuantitativo, reduciendo la ambigüedad y acelerando la lectura del estado de la app.
  \item \textbf{Verde \texttt{\#3D9A6F} como color primario}: Alineación semántica con sostenibilidad, diferenciado del verde genérico de Bootstrap.
  \item \textbf{Anillo SVG con gradiente}: Feedback inmediato y de alta fidelidad visual del progreso, con animación suave acelerada por hardware.
  \item \textbf{Estado centralizado en JavaScript}: Garantiza consistencia entre pestañas sin frameworks, manteniendo el código mantenible para equipos universitarios.
  \item \textbf{Bottom Sheet para donación}: Maximiza el alcance del pulgar y preserva el contexto visual del proyecto.
  \item \textbf{Sistema de microinteracciones}: Cada gesto del usuario recibe una respuesta física-visual inmediata, cerrando el bucle de feedback y reforzando el comportamiento.
\end{enumerate}

El resultado es una interfaz que \textbf{no necesita instrucciones}: el usuario comprende el modelo de Nativo en menos de 30 segundos de interacción, lo que es el indicador más preciso de éxito en IPO.



\newpage
% ============================================================
\section*{Referencias}
% ============================================================
\addcontentsline{toc}{section}{Referencias}

\begin{itemize}[leftmargin=*]

  \item Nielsen, J. (1994). \textit{10 Usability Heuristics for User Interface Design}. Nielsen Norman Group.
  \item ISO 9241-210:2010. \textit{Human-centred design for interactive systems}. ISO.
  \item Materiales de teoría de la asignatura Interacción Persona-Ordenador. (2025). \textit{Apuntes y diapositivas docentes del curso}.
  \item Statista. (2024). \textit{Mobile fitness app usage worldwide}. Statista Research Department.
  \item Google Fonts. (2024). \textit{Sora Font Family}. fonts.google.com.
  \item JetBrains. (2024). \textit{JetBrains Mono: A free and open source typeface}. jetbrains.com.
\end{itemize}

\end{document}
